\scp{Proto-Aru}{06-01}
\citet{Nivens1993,Nivens2017} posits the following regular sound changes
between PMP and Proto-Aru: *p > **ɸ,
*b > **p, *d > **r,
*j/*R > **R, *z/*y > **y,
and *q/*h > {\0}, with other consonants usually remaining unchanged.
Proto-Aru **ɸ was either a bilabial fricative [ɸ] or labiodental fricative [f].

In addition to the consonants which are reflexes of PMP proto-phonemes,
\citet{Nivens2017} also reconstructs three voiced stops for which
we cannot identify a consistent
and regular Austronesian source.
Our interpretation of the Proto-Aru phoneme inventory differs
slightly from that proposed by \citet{Nivens1993,Nivens2017}, in that we take the Proto-Aru voiced stops to be prenasalised.
The possible values of the Proto-Aru consonants which have changed
from PMP, as well as those which have no consistent PMP source are given
in \trf{06-02} along with their reflexes in selected languages.
The third line gives the phonetic values assigned to them by
\citet{Nivens1993,Nivens2017}, while the fourth line gives the interpretation which we
take in this chapter.

\begin{table}\tcp{Selected Proto-Aru consonants with selected reflexes}{06-02}
%\begin{adjustbox}{width=\textwidth}
	\begin{threeparttable} \st{4pt}
	\begin{tabular}{lllllllllll}\lsptoprule
PMP	&	\hr*p	&	\hr*b	&	\hr*d	&	*j/*R	&	*z/*y	&	\mc{2}{l}{(*mb)}			&		&		&	(*g)	\\	
Proto-Aru	&	**ɸ	&	**p	&	**r	&	**R	&	**y	&	\mc{2}{l}{**mb\su{‡}}			&	\mc{2}{l}{**nd}			&	**ŋg	\\	
Nivens	&	\hr[ɸ]/\hr[f]	&	\hr[p]	&	\hr[r]	&	\hr[ʀ]\su{†}	&	\hr[y]	&	\mc{2}{l}{\hr[b]}			&	\mc{2}{l}{\hr[d]}			&	\hr[g]	\\	
this chapter	&	\hr[ɸ]	&	\hr[p]	&	\hr[r]	&	\hr[ʁ]	&	\hr[y]	&	\mc{2}{l}{\hr[mb]}			&	\mc{2}{l}{\hr[nd]}			&	\hr[ŋg]	\\	
	&		&		&		&		&		&		&	/{\gap\#}	&		&	/{\gap\#}	&		\\	\midrule
Kompane	&	{\hr\hr}h	&	{\hr\hr}ɸ	&	{\hr\hr}r	&	{\hr\hr}h	&	{\hr\hr}ʤ	&	{\hr\hr}b 	&	-m	&	{\hr\hr}d	&	-n	&	{\hr\hr}ŋ	\\	
\tsc{Kola}	&	{\hr\hr}f	&	{\hr\hr}ɸʷ	&	{\hr\hr}r	&	{\hr\hr}h	&	{\hr\hr}y	&	{\hr\hr}b	&	-m	&	{\hr\hr}d	&	-n	&	{\hr\hr}g	\\	
Lorang	&	{\hr\hr}w	&	{\hr\hr}ɸ	&	{\hr\hr}r	&	{\hr\hr}r	&	{\hr\hr}ʤ	&	{\hr\hr}b	&	-m	&	{\hr\hr}d	&	-n	&	{\hr\hr}ŋ	\\	
North Dobel	&	{\hr\hr}w/y/{\0}	&	{\hr\hr}ɸ	&	{\hr\hr}r	&	{\hr\hr}r	&	{\hr\hr}s	&	{\hr\hr}b	&	-m	&	{\hr\hr}d	&	-n	&	{\hr\hr}ŋ	\\	
Barakai	&	{\hr\hr}{\0}	&	{\hr\hr}w	&	{\hr\hr}r	&	{\hr\hr}r	&	{\hr\hr}ʤ	&	{\hr\hr}b	&	-m	&	{\hr\hr}d	&	-n	&	{\hr\hr}g	\\	
SW. Tarangan	&	**w > {\0}	&	{\hr\hr}ɸ	&	{\hr\hr}r	&	{\hr\hr}r	&	{\hr\hr}ʤ	&	{\hr\hr}b	&	-m	&	{\hr\hr}d	&	-n	&	{\hr\hr}ŋ	\\	
		\lspbottomrule
	\end{tabular}
		\begin{tablenotes}\tnsize
			\item[†]	{\citet{Nivens1993} states that Proto-Aru **R was “possibly a uvular trill”.}
			\item[‡]	{\citet{Nivens1993,Nivens2017} transcribes the Proto-Aru voiced plosives as *b, *d,
								and *g in accordance with his interpretation of their phonetics.
								We have retranscribed them in accordance with our interpretation.}
		\end{tablenotes}
	\end{threeparttable}
%\end{adjustbox}
\end{table}

Evidence that the Proto-Aru voiced stops were prenasalised comes
from three main sources.
Firstly, while **mb and **nd are reflected as plain voiced stops
word initially and medially, they often (though not universally) change to nasals
when they become word final due to loss of a final vowel.
Two examples are given in reflexes of **ndumbu ‘six’
and **undu ‘blood’ in \trf{06-01},
with three additional examples given in \trf{06-03}.
If Proto-Aru
**mb and **nd were prenasalised this change is simply loss of the
plosive component word finally.
Thus, for instance, **undu > **und > Kola \it{un} ‘blood’.

Secondly, Proto-Aru **ŋg > \it{ŋ} in all languages except Kola,
Barakai, and some varieties of Dobel.
One example is Proto-Aru **ŋgamáR ‘sky’ given in \trf{06-01},
with three additional examples given in \trf{06-03}.
Again, if **ŋg were prenasalised [ŋg] this is simply loss of the plosive
element in languages with **ŋg > \it{ŋ}
and loss of prenasalisation in languages with **ŋg > \it{g}.

The final piece of evidence that Proto-Aru **mb
and **nd were prenasalised comes from the fact that there are
several reconstructions in which they are derived from PMP nasal-plosive
clusters, or correspond to such a cluster in other languages.
Examples include: *ampu > **(y)ambu ‘grandparent,
grandchild’, *Rambia ‘sago palm’ > **R[ae]mbi(y)a ‘sago starch’,
*paŋdan > **ɸandan ‘pandanus’,
and Proto-South Aru **wandan ‘Banda island’.

\begin{table}\tcp{Proto-Aru prenasalised stops}{06-03}
\begin{adjustbox}{width=\textwidth}
	\begin{threeparttable}\st{2pt}
	\begin{tabular}{lllllllll}\lsptoprule
local gl.	&	wound	&	leech	&	stone	&	nose, face	&	rattan	&	rotten\\
Proto-Aru	&	**tu\tbr{mb}u	&	**i\tbr{mb}i	&	**ku\tbr{mb}u	&	**lu\tbr{ŋg}a	&	**wa\tbr{ŋg}uR	&	**ya\tbr{ŋg}al-\\ \midrule
Ujir	&	{\hr\hr}tu\tbr{b}u	&	{\hr\hr}i\tbr{b}o	&		&	{\hr\hr}lu\tbr{ŋ}i	&	{\hr\hr}wa\tbr{ŋ}i	&	\\
Kompane	&	{\hr\hr}su\tbr{b}u	&	{\hr\hr}i\tbr{b}i	&	{\hr\hr}ku\tbr{b}u	&	{\hr\hr}nu\tbr{ŋ}ay	&	{\hr\hr}gwa\tbr{ŋ}u	&	\\
East Kola	&	{\hr\hr}su\tbr{m}	&	{\hr\hr}i\tbr{m}	&		&	{\hr\hr}lu\tbr{g}i	&	{\hr\hr}wa\tbr{g}uh	&	{\hr\hr}ya\tbr{g}il\\
Lola	&		&	{\hr\hr}yi\tbr{b}i	&	{\hr\hr}ʔu\tbr{b}u	&	{\hr\hr}lu\tbr{ŋ}i	&	{\hr\hr}wa\tbr{ŋ}	&	\\
Lorang	&	{\hr\hr}u\tbr{b}u	&	{\hr\hr}i\tbr{b}i	&		&	{\hr\hr}lu\tbr{ŋ}ay	&		&	{\hr\hr}sa\tbr{ŋ}ír\\
North Dobel	&		&	{\hr\hr}yi\tbr{b}i	&		&	{\hr\hr}lu\tbr{ŋ}ay-m	&	{\hr\hr}kwa\tbr{ŋ}ur	&	{\hr\hr}sa\tbr{ŋ}al-di\\
Dobel (Warjukur)	&		&		&	{\hr\hr}ʔu\tbr{b}u	&	{\hr\hr}lu\tbr{gw}ay	&	{\hr\hr}kwa\tbr{g}ur	&	{\hr\hr}sa\tbr{g}il\\
Koba	&		&	{\hr\hr}i\tbr{b}i	&		&	{\hr\hr}nu\tbr{ŋ}ay	&	{\hr\hr}ʍa\tbr{ŋ}ar	&	{\hr\hr}sa\tbr{ŋ}il\\
Barakai	&	{\hr\hr}so\tbr{m}	&	{\hr\hr}i\tbr{b}	&		&	{\hr\hr}nu\tbr{g}e	&	{\hr\hr}go\tbr{g}or	&	{\hr\hr}ʤá\tbr{g}el\\
Karey	&		&	{\hr\hr}gí\tbr{b}ado	&		&	{\hr\hr}no\tbr{ŋ}a\su{†}	&	{\hr\hr}gwa\tbr{ŋ}ur	&	{\hr\hr}ʤá\tbr{ŋ}el\\
Batuley	&	{\hr\hr}su\tbr{m}	&		&	{\hr\hr}ku\tbr{m}	&	{\hr\hr}nu\tbr{ŋ}ei	&	{\hr\hr}gwa\tbr{ŋ}or	&	{\hr\hr}ʤá\tbr{ŋ}el\\
Mariri	&	{\hr\hr}su\tbr{m}	&		&	{\hr\hr}ko\tbr{m}	&	{\hr\hr}dʊ\tbr{ŋ}ay	&	{\hr\hr}gwa\tbr{ŋ}or	&	{\hr\hr}ʤá\tbr{ŋ}el\\
Manombai	&	{\hr\hr}su\tbr{b}u	&	{\hr\hr}ibi	&		&	{\hr\hr}lu\tbr{ŋ}ay	&		&	{\hr\hr}ʤa\tbr{ŋ}ali\\
NE. Tarangan	&		&	{\hr\hr}i\tbr{b}a	&		&	{\hr\hr}lô\tbr{ŋ}ay\su{‡}	&	{\hr\hr}gwa\tbr{ŋ}ar	&	\\
SE. Tarangan	&		&	{\hr\hr}i\tbr{b}a	&		&	{\hr\hr}lô\tbr{ŋ}ay\su{‡}	&	{\hr\hr}gwa\tbr{ŋ}ar	&	\\
NW. Tarangan	&		&	{\hr\hr}ŋi\tbr{b}ar	&		&	{\hr\hr}lu\tbr{ŋ}ay	&	{\hr\hr}ga\tbr{ŋ}ar	&	{\hr\hr}ʤa\tbr{ŋ}il\\
SW. Tarangan	&	{\hr\hr}su\tbr{m}	&	{\hr\hr}i\tbr{m}	&		&	{\hr\hr}lu\tbr{ŋ}ay	&	{\hr\hr}ga\tbr{ŋ}ar	&	{\hr\hr}ʤa\tbr{ŋ}il\\
		\lspbottomrule
	\end{tabular}
		\begin{tablenotes}\tnsize
			\item[†]	{Karey \it{noŋa} is ‘forehead’.}
			\item[‡]	{NE. and SE. Tarangan \it{lônay} ‘nose, face’ is from Gomar Meti
								and Dosi Mar villages respectively.}
		\end{tablenotes}
	\end{threeparttable}
\end{adjustbox}
\end{table}

The principal evidence for identifying the \tsc{Aru} languages
as an innovation-defined subgroup within Malayo-Polynesian is
merger of *j/*R, merger of *z/*y, and sporadic glide insertion
before certain vowel initial words.
The changes *p > **ɸ
and *b > **p can be added as supporting evidence.
